\section{The Gifts Prepared and the Church Implicated}

\massunit{14}{Offertory antiphon and offering of the host}{Proper slot plus fixed prayer}{\latin{Ordo Missae}, p.~220}

\textbf{Literal and rubrical action.} The proper Offertory antiphon supplies the day's scriptural voice. The celebrant uncovers the chalice, takes the paten with the bread, raises it, and says \latin{Suscipe, sancte Pater}; he then signs the paten and places the host on the corporal. In Solemn Mass the subdeacon brings the covered chalice and paten from the credence, the deacon presents the paten and assists according to his office, and the schola sings the proper Offertory. This differentiated 1962 preparation must not be confused with the earlier Roman people's procession of gifts, although that public action belongs to the Offertory chant's history. The printed prayer calls the still-unconsecrated bread an ``immaculate host'' and speaks of sins, the living and dead, and salvation; its grammar is anticipatory and must be read toward the Canon, not as a claim that transubstantiation has already occurred.

\textbf{Scripture and history.} Christians bring created goods because all creation belongs to God and is received with thanksgiving (Gen.~14:18--20; 1~Chr.~29:10--14; 1~Tim.~4:4--5). Justin describes bread, wine, and water being brought after the prayers and kiss (\emph{First Apology} 65, 67); Irenaeus interprets the Church's new-covenant offering as a confession that the Creator gives firstfruits (\emph{Against Heresies} IV.17.5; IV.18.1--6). The Offertory chant accompanied a once-prominent procession of gifts in the Roman order. \latin{Suscipe, sancte Pater}, however, belongs to the later medieval private offertory complex and should not be attributed verbatim to Justin or Irenaeus.

\textbf{Doctrine, symbolism, and relation.} Bread condenses creation and human labor into matter apt for Christ's institution. The priest's singular unworthiness and wide intercession connect personal ministry to the whole Church, but the prayer does not make unconsecrated bread the sacramental Victim. Trent distinguishes the preparation of gifts from the moment in which the same Christ offered on Calvary is contained and unbloodily offered (Session XXII, chs.~1--2). Symbolically, hearing has become handling: the faith confessed in unit~13 now takes created form and accepts implication in sacrifice. Unit~14 begins a preparation whose divine dependence will become explicit in units~16 and 20.

\begin{massbox}{Proleptic language}
Traditional liturgical speech may name a thing from the end toward which the rite orders it. ``Host'' here is proleptic, not a denial of the dogmatic distinction between bread before consecration and Christ's Body after it.
\end{massbox}

\massunit{15}{Water and wine; offering of the chalice}{Fixed}{\latin{Ordo Missae}, pp.~220--221}

\textbf{Literal and rubrical action.} Wine is poured into the chalice and a small quantity of water is mixed with it. The priest blesses the water except in Masses for the dead, praying through the mystery of the water and wine that human beings may share Christ's divinity, he who shared humanity. He then raises the chalice and says \latin{Offerimus tibi}, places it on the corporal, and covers it. The singular host-prayer and plural chalice-prayer should not be made into a speculative opposition: both belong to one preparation by priest and Church.

\textbf{Scripture and patristic witness.} Christ took the cup of wine at the Supper (Matt.~26:27--29; 1~Cor.~11:25), and blood and water flow from his pierced side (John~19:34). Cyprian gives the decisive early Western dossier: water alone or wine alone fails the dominical and ecclesial sign; wine signifies Christ's Blood, water the people joined to him, and their mixture cannot afterward be separated (\emph{Epistle} 63.2, 9, 13). His exegesis is a received patristic interpretation, while the historical practice of mixed wine is independently attested by Justin (\emph{First Apology} 65).

\textbf{Doctrine, symbolism, and relation.} The mixing prayer draws on 2~Peter~1:4 and the Incarnation to frame sacramental participation: creatures do not become divine by nature, and humanity is not absorbed into God; grace gives a created participation in divine life through Christ. The water does not represent a second victim or add a human cause alongside Christ. Trent explicitly defends mixing water with wine, citing both Christ's example and the blood-and-water mystery (Session XXII, ch.~7). Unit~15 expands the host's created sign into covenantal and ecclesial communion. It leads to unit~16's admission that even rightly prepared matter cannot sanctify itself.

\massunit{16}{\latin{In spiritu humilitatis} and \latin{Veni, sanctificator}}{Fixed}{\latin{Ordo Missae}, p.~221}

\textbf{Literal and rubrical action.} Bowed, the celebrant prays to be received ``in a spirit of humility'' and asks that the sacrifice be made pleasing. He then raises his eyes, extends and lifts his hands, invokes the eternal and almighty God as Sanctifier, signs both offerings, and asks blessing. These prayers still belong to preparation. Neither the sign nor the imperative \latin{Veni} is the sacramental form of consecration.

\textbf{Scriptural and historical horizon.} \latin{In spiritu humilitatis} takes its grammar from the prayer of Azariah in the furnace, where Israel without cultic resources asks that contrition be received like sacrificial abundance (Dan.~3:39--40 in the Vulgate/Greek additions). The prayer's insertion into the medieval offertory makes the Church's material preparation morally accountable. \latin{Veni, sanctificator} belongs to the same medieval Roman-Frankish development. Because its address is ``eternal, almighty God,'' it should not be presented without qualification as an explicit Eastern-form epiclesis to the Holy Spirit.

\textbf{Doctrine, symbolism, and relation.} Sacrifice is acceptable through God's initiative, not because bread, wine, ceremonial precision, or ministerial intensity compel him. Aquinas locates sacramental conversion in Christ's words pronounced by the priest in Christ's person (\emph{ST} III, q.~78, aa.~1--4; q.~82, a.~1). The preparatory invocation is therefore real prayer for blessing and acceptance without competing with the form to come. The pairing also corrects two errors at once: matter without humble self-offering becomes externalism; interior sentiment without the received gifts and rite becomes disembodied. Unit~16 turns offering into supplication and prepares the upward sign of incense in unit~17.

\massunit{17}{Incensation of gifts, cross, altar, ministers, and people}{Conditional solemn action}{\latin{Ordo Missae}, p.~221; \latin{Ritus servandus} VII}

\textbf{Literal and rubrical action.} In Solemn Mass the celebrant blesses incense, incenses the offerings with prescribed crosses and circles, then cross and altar; the deacon receives the thurible and incenses celebrant, ministers, clergy, and people according to rank. The accompanying prayers ask that incense blessed by God rise as sweetness, that the Lord enkindle charity, and that the prayer ascend like incense. A Low Mass simply lacks this complex. Its absence does not reduce sacramental efficacy, and its presence does not confer identical liturgical office on everyone incensed.

\textbf{Scriptural and historical horizon.} Psalm~140:2 gives the explicit comparison between prayer and incense; Malachi~1:11 and Revelation~8:3--4 locate fragrant offering within worship extending through earth and heaven. Incense is present in early Roman ceremonial, particularly around entrance and Gospel, while the elaborate offertory incensation and its private prayers matured in the medieval rite. The loops and signs should therefore be described from the 1962 rubric before any spiritual allegory is attached to them.

\textbf{Patristic, doctrinal, and relational force.} Augustine reads the ``sacrifice of praise'' as the worship of the redeemed city offered through Christ the High Priest (\emph{City of God} X.6). Incense makes ascent sensible, but smoke is neither the Holy Spirit nor the substance of the Church's oblation. Its ordered route gathers gifts, altar, ministerial body, and faithful into relation to the one sacrifice without erasing hierarchy. The gifts are incensed first because they are ordered toward consecration; persons are incensed because they belong to the worshipping body. Unit~17 exteriorizes the plea of unit~16 and leads naturally to unit~18, where the minister asks that the one who handles holy things be cleansed.

\massunit{18}{Lavabo}{Fixed; doxology conditionally omitted}{\latin{Ordo Missae}, pp.~221--222}

\textbf{Literal and rubrical action.} At the Epistle side the priest washes the tips of thumb and forefinger while reciting Psalm~25:6--12, then returns to the middle. The doxology is omitted in Masses for the dead and during the specified Passiontide period. The act retains a practical relation to hands that have received offerings and incense, but the psalm makes its ritual meaning explicitly moral and cultic: innocence, altar, thanksgiving, and love of God's dwelling.

\textbf{Scripture and patristic witness.} The psalm's ``washing among the innocent'' differs from Pilate's evasion of culpability (Matt.~27:24): the celebrant asks purification so as to remain answerable before God. Cyril of Jerusalem explains the pre-anaphoral washing as a symbol that ministers must be pure from sinful deeds, insisting it is not ordinary washing for bodily dirt (\emph{Mystagogical Catechesis} V.2). Cyril describes Jerusalem, not the medieval Roman Offertory, but his exact mystagogical distinction responsibly illuminates the shared sign.

\textbf{Doctrine, symbolism, and relation.} Exterior washing cannot cleanse mortal guilt by itself; it petitions and teaches. \emph{Mediator Dei} insists that interior and exterior worship belong together and condemns both empty formalism and contempt for bodily rites (23--24). The hands signify agency: the priest truly acts, yet needs grace to act worthily. Psalm~25 also shifts from singular plea to assembly---the psalmist will bless God ``in the churches''---so purification is ordered to ecclesial praise. Unit~18 does not restart the penitential rite; it intensifies readiness at the point where prepared gifts are about to be commended in the name of the Trinity.

\massunit{19}{\latin{Suscipe, sancta Trinitas} and \latin{Orate, fratres}}{Fixed}{\latin{Ordo Missae}, p.~222}

\textbf{Literal and rubrical action.} Bowed at the altar, the celebrant commends the oblation to the Holy Trinity in memory of Christ's Passion, Resurrection, and Ascension and in honor of Mary and the saints. He kisses the altar, turns toward the people, and asks them to pray that ``my sacrifice and yours'' be acceptable; the server or ministers answer that the Lord receive the sacrifice for divine praise and glory, the offerers' good, and that of the whole Church. The priest then answers Amen quietly and turns to the Secret.

\textbf{Scriptural and historical horizon.} The Paschal triad rests on the apostolic proclamation that the one who died was raised and exalted (1~Cor.~15:3--5; Phil.~2:6--11; Heb.~9:24). Hebrews~13:15--16 joins Christ-mediated praise with beneficence and shared gifts. The exact \latin{Suscipe} is a medieval synthesis; the invitation \latin{Orate, fratres} and its response also reached their settled form through medieval development. The saint clause means that honor given them returns to God's glory and their intercession aids the faithful; it does not make saints co-victims or rival recipients of divine sacrifice.

\textbf{Doctrine, symbolism, and relational force.} ``My and yours'' distinguishes while uniting. The ordained priest alone acts sacramentally in the person of Christ to consecrate; the faithful offer the divine Victim through his hands and offer themselves in union with Christ according to the common priesthood (\emph{Mediator Dei} 80--104; \emph{Lumen gentium} 10--11). Neither clerical isolation nor an undifferentiated co-consecration fits the dialogue. The response supplies the fourfold horizon of worship: God's glory, sanctification of offerers, ecclesial good, and universality. Unit~19 makes explicit what the plural chalice prayer anticipated; unit~20 now gives the Church's proper, day-specific prayer over these gifts.

\begin{studybox}{The Paschal measure of every offering}
The gifts are not religious tribute added to Christ's work. They become ecclesial sacrifice only within the memorial of his Passion, Resurrection, and Ascension and through the consecratory action that follows. Self-offering is Christian when caught up into his prior oblation.
\end{studybox}

\massunit{20}{Secret}{Proper slot}{\latin{Ordo Missae}, p.~222}

\textbf{Literal and rubrical action.} With hands extended, the celebrant says the appointed Secret or Secrets quietly. The conclusion is joined to the opening of the Preface: at its final words he says or sings \latin{Per omnia saecula saeculorum} aloud, the people answer Amen, and the dialogue begins. The prayer is a stable functional slot with variable text, normally asking acceptance, sanctification, or fruit in relation to the gifts and feast. Its quiet delivery belongs to the received 1962 mode and should not be explained as though inaudibility caused efficacy.

\textbf{History and theological grammar.} The Secret descends from the ancient Roman prayer over the offerings preserved in sacramentaries; unlike most private offertory prayers, it belongs to the old presidential prayer structure. The origin of the name is disputed: derivations from ``set apart'' and from quiet recitation have both been proposed, so a single etymology should not be treated as settled history. Aquinas places the prayer asking that the people's oblation be acceptable immediately before the Preface in his account of the Mass's ordered words (\emph{ST} III, q.~83, a.~4).

\textbf{Symbolic and relational force.} The Secret seals preparation but does not form a barrier between Church and priest. Its plural prayer presupposes the public exchange of unit~19, and the people's Amen ratifies its conclusion. Quiet concentrates attention on the altar while the audible ending reopens the action into common thanksgiving. The proper text also prevents the Offertory from becoming a timeless generic theory: each celebration asks that these gifts bear fruit according to the mystery commemorated. Unit~20 is therefore a hinge, not a terminal prayer. Petition rises into the Preface because acceptance of the gifts can only be thanksgiving within Christ's own self-offering.

\begin{massbox}{Proper-slot caution}
No universal exposition may attribute to every Secret a petition found only in one formulary. The complete text, grammar, and conclusion of the day's prayer govern its local theological claims.
\end{massbox}
